\documentclass[12pt,a4paper,oneside]{book}
\usepackage[utf8]{inputenc}
\usepackage[T1]{fontenc}
\usepackage[english,french]{babel}
\usepackage{graphicx}
\usepackage{geometry}
\usepackage[table]{xcolor}
\usepackage{colortbl}
\usepackage{titlesec}
\usepackage{fancyhdr}
\usepackage{hyperref}
\usepackage{xcolor}
\usepackage{listings}
\usepackage{multicol}
\usepackage{amsmath}
\usepackage{amssymb}
\usepackage{booktabs}
\usepackage{longtable}
\usepackage{parskip}
\usepackage{enumitem}
\usepackage{fontawesome5}
\usepackage{url}

\definecolor{electronblue}{RGB}{0, 120, 212}
\definecolor{electronpurple}{RGB}{136, 23, 152}
\definecolor{electroncyan}{RGB}{0, 190, 255}
\definecolor{electrongold}{RGB}{255, 185, 0}
\definecolor{darkbg}{RGB}{15, 23, 42}
\definecolor{lightbg}{RGB}{249, 250, 251}
\definecolor{softgray}{RGB}{248, 249, 250}
\definecolor{mediumgray}{RGB}{107, 114, 128}
\definecolor{success}{RGB}{34, 197, 94}
\definecolor{warning}{RGB}{251, 191, 36}

\geometry{top=2.5cm,bottom=2.5cm,left=2.5cm,right=2.5cm}

\hypersetup{
    colorlinks=true,
    linkcolor=electronblue,
    filecolor=electronpurple,
    urlcolor=electroncyan,
    pdfauthor={E-Graphisme by ELECTRON},
    pdftitle={Documentation Enterprise - Volume 6},
    pdfsubject={Architecture Technique & Infrastructure}
}

\pagestyle{fancy}
\fancyhf{}
\fancyhead[LE,RO]{\thepage}
\fancyhead[RE,LO]{\leftmark}
\fancyfoot[CE,CO]{\thepage}

\titleformat{\chapter}[display]
{\normalfont\Huge\bfseries\color{electronblue}}
{\chaptertitlename\ \thechapter}{20pt}{\Huge}
\titleformat{\section}
{\normalfont\Large\bfseries\color{electronblue}}
{\thesection}{1em}{}
\titleformat{\subsection}
{\normalfont\large\bfseries\color{electronpurple}}
{\thesubsection}{1em}{}

\begin{document}

\begin{titlepage}
\thispagestyle{empty}
\begin{center}
\vspace*{2cm}
{\Huge\bfseries\color{electronblue}E-GRAPHISME}\\[0.3cm]
{\LARGE\bfseries\color{electronpurple}by ELECTRON}\\[2cm]
\vspace{2cm}
{\huge\bfseries\color{darkbg}ARCHITECTURE TECHNIQUE}\\[0.3cm]
{\huge\bfseries\color{darkbg}\& INFRASTRUCTURE}\\[1cm]
{\Large\bfseries\color{electroncyan}Volume 6}\\[3cm]
\vspace{2cm}
{\Large\color{mediumgray}Architecture • Performance • Sécurité • Évolutivité}\\[1cm]
{\Large\color{mediumgray}Là où la créativité rencontre la technologie}\\[3cm]
{\Large\bfseries\color{electrongold}Version 2026}\\[2cm]
\vfill
\end{center}
\end{titlepage}

\frontmatter

\chapter*{Préface}
\addcontentsline{toc}{chapter}{Préface}
Ce volume présente l'architecture technique et l'infrastructure de la plateforme E-Graphisme by ELECTRON.

\bigskip
\textbf{Classification :} Document Interne - Partageable\\
\textbf{Volume :} 6\\
\textbf{Version :} 2026

\tableofcontents

\mainmatter

\chapter{VISION}

\section{Principes}
L'architecture de E-Graphisme by ELECTRON est conçue comme une plateforme modulaire, évolutive et hautement disponible.

Chaque composant peut évoluer indépendamment sans interrompre le fonctionnement global.

\section{Modes de Fonctionnement}
Le système est pensé pour fonctionner :
\begin{itemize}
\item entièrement en local ;
\item sur un serveur privé ;
\item dans un cloud privé ;
\item ou en mode hybride.
\end{itemize}

\chapter{ARCHITECTURE GLOBALE}

\section{Couches de l'Architecture}
L'application est divisée en plusieurs couches :

\begin{enumerate}
\item Utilisateurs
\item Frontend Premium
\item API Gateway
\item Services Métiers
\item AI Orchestrator
\item Agents IA
\item Ollama
\item Base de données
\item Stockage
\item Sauvegardes
\end{enumerate}

Chaque couche est indépendante et communique via des API internes.

\chapter{INFRASTRUCTURE LOCALE}

\section{Composants}
Le serveur local comprend :
\begin{multicols}{2}
\begin{itemize}
\item Docker Desktop
\item Ollama
\item OpenHands
\item PostgreSQL
\item Redis
\item Firebase (optionnel)
\item Serveur Node.js
\item Serveur React
\end{itemize}
\end{multicols}

Tous les services sont conteneurisés.

\chapter{DOCKER ENTERPRISE}

\section{Isolation des Services}
Chaque module est isolé dans un conteneur dédié.

\section{Conteneurs Principaux}
\begin{center}
\begin{tabular}{|c|p{10cm}|}
\hline
\rowcolor{electronblue}
\color{white}\textbf{Conteneur} & \color{white}\textbf{Rôle} \\
\hline
Frontend & Application web React/Next.js \\
\hline
Backend & API Node.js/FastAPI \\
\hline
API Gateway & Point d'entrée sécurisé \\
\hline
PostgreSQL & Base de données relationnelle \\
\hline
Redis & Cache et sessions \\
\hline
Ollama & Moteur IA \\
\hline
OpenHands & Automatisation \\
\hline
AI Orchestrator & Coordination des agents \\
\hline
Print Service & Impression \\
\hline
Marketplace & Boutique \\
\hline
Monitoring & Supervision \\
\hline
\end{tabular}
\end{center}

Chaque conteneur possède :
\begin{itemize}
\item logs ;
\item sauvegardes ;
\item surveillance ;
\item redémarrage automatique.
\end{itemize}

\chapter{OLLAMA AI SERVER}

\section{Présentation}
Ollama constitue le moteur IA principal. Il fournit les modèles de langage utilisés par les agents.

\section{Gestion des Modèles}
Chaque agent peut être associé à un modèle spécifique selon sa spécialité.

Le système doit permettre :
\begin{itemize}
\item la gestion des modèles installés ;
\item la sélection du modèle par agent ;
\item le suivi des performances ;
\item la mise à jour des modèles.
\end{itemize}

\chapter{OPENHANDS}

\section{Rôle}
OpenHands agit comme l'assistant de développement logiciel.

\section{Interventions}
Il intervient pour :
\begin{itemize}
\item générer du code ;
\item corriger des erreurs ;
\item proposer des améliorations ;
\item maintenir la documentation ;
\item assister les développeurs.
\end{itemize}

Toutes les actions restent sous validation humaine avant intégration.

\chapter{FRONTEND}

\section{Technologies}
\begin{center}
\begin{tabular}{|c|c|}
\hline
\rowcolor{electronpurple}
\color{white}\textbf{Technologie} & \color{white}\textbf{Usage} \\
\hline
React & Bibliothèque UI \\
\hline
Next.js & Framework SSR/SSG \\
\hline
TypeScript & Typage statique \\
\hline
Tailwind CSS & Framework CSS \\
\hline
Framer Motion & Animations \\
\hline
Three.js & Graphiques 3D \\
\hline
\end{tabular}
\end{center}

\section{Principes}
\begin{itemize}
\item responsive ;
\item accessibilité ;
\item animations fluides ;
\item composants réutilisables ;
\item optimisation SEO.
\end{itemize}

\chapter{BACKEND}

\section{Modules}
Le Backend est organisé en modules :

\begin{center}
\begin{tabular}{|c|p{10cm}|}
\hline
\rowcolor{electroncyan}
\color{white}\textbf{Module} & \color{white}\textbf{Description} \\
\hline
Authentification & Connexion et sécurité \\
\hline
Utilisateurs & Gestion des comptes \\
\hline
Clients & CRM \\
\hline
Projets & Suivi projet \\
\hline
Marketplace & Boutique en ligne \\
\hline
Boutique & Produits \\
\hline
Paiements & Transactions \\
\hline
IA & Intelligence artificielle \\
\hline
Fichiers & Stockage \\
\hline
Notifications & Alertes \\
\hline
Rapports & Données analytiques \\
\hline
\end{tabular}
\end{center}

Chaque module possède ses propres services et contrôleurs.

\chapter{API GATEWAY}

\section{Fonctions}
Toutes les communications passent par une passerelle API.

\begin{center}
\begin{tabular}{|c|p{10cm}|}
\hline
\rowcolor{electronblue}
\color{white}\textbf{Fonction} & \color{white}\textbf{Description} \\
\hline
Authentification & Vérification des tokens \\
\hline
Routage & Distribution des requêtes \\
\hline
Limitation de débit & Rate limiting \\
\hline
Journalisation & Logging complet \\
\hline
Sécurité & Protection \\
\hline
\end{tabular}
\end{center}

\chapter{BASES DE DONNÉES}

\section{Types de Données}
Le système distingue plusieurs types de données.

\section{PostgreSQL}
\begin{itemize}
\item utilisateurs ;
\item commandes ;
\item factures ;
\item projets ;
\item CRM.
\end{itemize}

\section{Redis}
\begin{itemize}
\item cache ;
\item sessions ;
\item files d'attente.
\end{itemize}

\section{Stockage Documentaire}
\begin{itemize}
\item images ;
\item vidéos ;
\item PDF ;
\item contrats ;
\item maquettes.
\end{itemize}

\chapter{GESTION DES FICHIERS}

\section{Catégories}
Les ressources sont classées par catégories :
\begin{multicols}{2}
\begin{itemize}
\item œuvres d'art ;
\item logos ;
\item vidéos ;
\item documents ;
\item livrables ;
\item modèles.
\end{itemize}
\end{multicols}

\section{Métadonnées}
Chaque fichier possède :
\begin{itemize}
\item un identifiant unique ;
\item un historique de versions ;
\item des métadonnées ;
\item des droits d'accès.
\end{itemize}

\chapter{SÉCURITÉ}

\section{Niveaux de Protection}
La sécurité repose sur plusieurs niveaux :
\begin{itemize}
\item authentification forte ;
\item contrôle d'accès par rôles ;
\item chiffrement des données sensibles ;
\item journal d'audit ;
\item sauvegardes régulières ;
\item surveillance des connexions.
\end{itemize}

Le Centre IA est accessible uniquement au Fondateur.

\chapter{MONITORING}

\section{Éléments Surveillés}
Le système surveille en continu :
\begin{multicols}{2}
\begin{itemize}
\item l'état des conteneurs Docker ;
\item les performances CPU et GPU ;
\item l'utilisation de la mémoire ;
\item les bases de données ;
\item les agents IA ;
\item les temps de réponse.
\end{itemize}
\end{multicols}

Les alertes critiques sont affichées sur le tableau de bord du Fondateur.

\chapter{SAUVEGARDES}

\section{Éléments Sauvegardés}
Les sauvegardes portent sur :
\begin{multicols}{2}
\begin{itemize}
\item bases de données ;
\item documents ;
\item paramètres ;
\item prompts des agents IA ;
\item journaux.
\end{itemize}
\end{multicols}

Des restaurations ciblées ou complètes peuvent être réalisées.

\chapter{DÉPLOIEMENT}

\section{Plateformes Cibles}
La plateforme peut être installée sur :
\begin{itemize}
\item un ordinateur personnel ;
\item un serveur local ;
\item une machine virtuelle ;
\item un serveur dédié ;
\item une infrastructure cloud.
\end{itemize}

Le déploiement est piloté par Docker Compose.

\chapter{ÉVOLUTIONS FUTURES}

\section{Modules Prévus}
L'architecture est conçue pour intégrer de nouveaux modules :
\begin{itemize}
\item application mobile ;
\item synchronisation multi-sites ;
\item fédération de plusieurs agences ;
\item marketplace internationale ;
\item nouveaux modèles IA ;
\item services cloud optionnels.
\end{itemize}

\appendix

\chapter{STRUCTURE DU PROJET}

\section{Organisation}
Le dépôt source sera organisé en grands espaces fonctionnels :

\section{Applications}
\begin{itemize}
\item site public ;
\item ERP ;
\item Marketplace ;
\item Studio IA ;
\item Portail Client ;
\item Portail Affilié ;
\item Centre IA.
\end{itemize}

\section{Services}
\begin{itemize}
\item API ;
\item moteur d'orchestration IA ;
\item impression ;
\item notifications ;
\item paiements ;
\item stockage.
\end{itemize}

\section{Bibliothèques}
\begin{itemize}
\item composants UI ;
\item thèmes ;
\item icônes ;
\item utilitaires ;
\item SDK internes.
\end{itemize}

\section{Configuration}
\begin{itemize}
\item Docker ;
\item variables d'environnement ;
\item sécurité ;
\item journalisation.
\end{itemize}

\section{Documentation}
\begin{itemize}
\item guides ;
\item API ;
\item procédures ;
\item architecture.
\end{itemize}

\section{Tests}
\begin{itemize}
\item unitaires ;
\item intégration ;
\item performances ;
\item sécurité.
\end{itemize}

Cette organisation facilite la maintenance, les mises à jour et l'ajout de nouveaux modules.

\chapter{Index}
\begin{enumerate}
\item Vision de l'Architecture
\item Architecture Globale
\item Infrastructure Locale
\item Docker Enterprise
\item Ollama AI Server
\item OpenHands
\item Frontend
\item Backend
\item API Gateway
\item Bases de données
\item Gestion des fichiers
\item Sécurité
\item Monitoring
\item Sauvegardes
\item Déploiement
\item Évolutions futures
\item Structure du Projet
\end{enumerate}

\end{document}
